\note{2}{Wed 27 Sep 2023 17:49}{Conceptual Summary}

\section{Topological Quantum Computing}
\subsection{Summary of important concepts}
\subsubsection{What is TQC}

Topological Quantum Computing performs quantum computation using the topological properties of anyons. 

Three main steps of TQC:
\begin{itemize}
	\item \textbf{State Preparation}. We do \textbf{pair production} and \textbf{fusion} to make several anyons in the system.
	\item \textbf{Braiding.} We swap the anyons around to edit the topological state.
	\item \textbf{Fusion}. We may fuse some anyons together and detect the emission. 
\end{itemize}

\subsubsection{Formalization}

Types of anyons and their interactions are described by the \textbf{Quantum Double}, a quasi-triangular Hopf algebra $DG$ over some group $G$. Specifically, the \textit{irreducible representations} of DG correspond to the types of anyons. 

Braiding of anyons, or the ``worldlines'' of anyons (they are actually ``ribbons'' instead of ``lines''), are described by the \textbf{Modular Tensor Category}. A worldline that starts with anyons $a$ and $b$ and ends with anyons $c$ and $d$ is an element in $\operatorname{Hom}(a \otimes b, c \otimes d)$. We regard ``no particle'' as one anyon type called the \textit{vacuum}, denoted $\mathbbm{1}$. Objects in MTC are finite linear superpositions of such worldlines.
The MTC gives us a theory describing what observation we get for arbitrary braiding and fusing of anyons.

\subsubsection{Operations to worldlines}

We can physically manipulate the system in multiple ways:
\begin{itemize}
	\item We can perform \textbf{Braiding} on the anyons, by swapping them around using some (technologies still under development). By doing this we alter the quantum state encoded as some object of the MTC.
	\item By forming \textbf{Fusing} on the anyons, we let two anyons react to produce other anyons.
	Fusing of anyons are given by the decomposition of product of irreps:
\[
	a \otimes b = \sum_{\gamma} N_{a b}^\gamma \gamma
.\] 

	It is possible to \textbf{pair produce} an anyon $a$ and its dual (antiparticle) $\overline{a} $ from the vacuum. In the MTC, we are composing arrows.

	\item As a special case of Fusing, we can \textbf{taking trace} of some object in the category, i.e. make a worldline that closes onto itself (some endomorphism in the MTC), the anyons will vanish and produce observable emissions. This allows us to extract data from the ``computer''.
	\item \textbf{Measurements} are in principle possible through inteferometer. This can collapse a superposition state to a ground state.
\end{itemize}

\subsubsection{Effect of braiding}

Braidings alter the phase and amplitude of topological states. The idea is to analyze a small set of operations that generate the whole MTC. Swap and Twist introduce a phase to the quantum state; Link gives an emission; Slide allow us to ``slide the worldlines around''. 

\begin{itemize}
	\item \textbf{Swap.} 
	\item \textbf{Twist.}
	\item \textbf{Link.}
	\item \textbf{Change basis (Slide).}
\end{itemize}
All those four operations can be explicitly calculated for a given MTC.

To perform computation, we need to do two things:
\begin{itemize}
	\item Construct a basic \textbf{computational object} from available anyons. This is usually a space $\operatorname{Hom} (a, b_1, \ldots, b_k)$. Those $b$ 's need not be distinct. The dimension of the hom space represents number of distinct ways (up to global phase) to start from a single anyon $a$ and produce the ordered collection of particles $b_1 \ldots b_k$. When $\operatorname{dim}  = 2$ this is a \textit{qubit}, when $\operatorname{dim}  = 3$ this is a \textit{qutrit}, etc. 
	\item \textbf{Quantum gates} are given by braidings on single or multiple qubits. A quantum gate is \textit{entangling} if it is not the tensor product of multiple single-qubit gates. When we put multiple qubits together, their total $\operatorname{Hom} $ space can be strictly larger than the logical space, tensor product of the qubits. We need our gate to be \textbf{leakage-free} in the sense that it preserves the logical space. To make it a strong quantum computer, we want to find a \textbf{universal set} of gates. The existence of leakage-free and universal set of gates depend on specific TQFTs, and is solved only for a handful of cases.
\end{itemize}

\subsection{Quantum Double}

\subsubsection{Construction}

\[
	DG = \mathbb{C}[G]^{*, \text{cop}} \otimes \mathbb{C}[G]
.\] 


\subsubsection{Representations of DG}

A homomorphism $\phi$ between representations $(\chi, V)$ and $(\chi', W)$ must satisfy the commutative diagram for all $x \in G$ :

% https://q.uiver.app/#q=WzAsNCxbMCwwLCJWIl0sWzMsMCwiViJdLFswLDIsIlciXSxbMywyLCJXIl0sWzAsMiwiXFxwaGkiXSxbMSwzLCJcXHBoaSJdLFswLDEsIlxcY2hpKHgpIiwxXSxbMiwzLCJcXGNoaScoeCkiLDFdXQ==
\[\begin{tikzcd}
	V &&& V \\
	\\
	W &&& W
	\arrow["\phi", from=1-1, to=3-1]
	\arrow["\phi", from=1-4, to=3-4]
	\arrow["{\chi(x)}"{description}, from=1-1, to=1-4]
	\arrow["{\chi'(x)}"{description}, from=3-1, to=3-4]
\end{tikzcd}\]

The comultiplication $\Delta$ allows us to multiply two representations:
\[
	(\chi \otimes  \chi') (x) := (\chi \otimes \chi')(\Delta(x))
.\] 

$\text{Rep}(DG)$ forms a monoidal category under the tensor product $\otimes $.

Each representation is a unique direct sum of irreps; in other words, $\text{Rep}(DG)$ is a semisimple category.

\subsubsection{Characterization of irreps}

Each irrep of $DG$ is characterized by a pair $(C, \chi)$ where $C$ is a conjugate class of $G$, and $\chi$ is an irrep of $Z(r)$ for some $r \in C$. (choice of $r$ is irrelevant)

Specific definition of $\rho_{(C, \chi)}$ :
\begin{align*}
	V_\rho &= \mathbb{C}^{\otimes C} \otimes V_\chi \\
	B_h . \left| c \right\rangle \otimes \left| v \right\rangle 
	&= \delta_{h,c} \left| c \right\rangle \otimes \left| v \right\rangle  \\
	A_g . \left| c \right\rangle  \otimes \left| v \right\rangle 
	&= \left| g c \overline{g}  \right\rangle \otimes \chi\left( \overline{q} _{g c \overline{g} } g q_c \right) \left| v \right\rangle  
.\end{align*}
where $q_c$ is chosen such that $q_c r q_c = c$.


\subsubsection{Correspondence between anyon type and irrep}



\subsection{Modular Tensor Category}


\subsection{Quantum bits and Quantum Gates}







